
\documentclass{article}
\usepackage{colortbl}
\usepackage{makecell}
\usepackage{multirow}
\usepackage{supertabular}

\begin{document}

\newcounter{utterance}

\centering \large Interaction Transcript for game `cladder', experiment `full\_v1.5\_default', episode 6328 with Qwen3.5{-}9B.
\vspace{24pt}

{ \footnotesize  \setcounter{utterance}{1}
\setlength{\tabcolsep}{0pt}
\begin{supertabular}{c@{$\;$}|p{.15\linewidth}@{}p{.15\linewidth}p{.15\linewidth}p{.15\linewidth}p{.15\linewidth}p{.15\linewidth}}
    \# & \multicolumn{2}{c}{Player} && \multicolumn{2}{c}{Game Master} \\
    \hline

    \theutterance \stepcounter{utterance}  
    & & & \multicolumn{4}{p{0.6\linewidth}}{
        \cellcolor[rgb]{0.9,0.9,0.9}{
            \makecell[{{p{\linewidth}}}]{
                \texttt{\tiny{[P1$\langle$GM]}}
                \texttt{You are an expert in causal inference. The following question is not a typical commonsense query, but rather a meticulously designed question created by a professor specializing in causal inference, intended to assess the students' mastery of the course content.} \\
\texttt{Imagine a self{-}contained, hypothetical world with only the following conditions, and without any unmentioned factors or causal relationships: Swoy has a direct effect on rixq and xevu. Rixq has a direct effect on zuph. Zuph has a direct effect on xevu. Swoy is unobserved. For those who are not rixq and are not zuph, the probability of xevu is 46\%. For those who are not rixq and are zuph, the probability of xevu is 41\%. For those who are rixq and are not zuph, the probability of xevu is 66\%. For those who are rixq and are zuph, the probability of xevu is 60\%. For those who are not rixq, the probability of zuph is 61\%. For those who are rixq, the probability of zuph is 34\%. For those who are rixq, would it be more likely to see xevu if the individual was not rixq?} \\
\texttt{Start your answer with 'Yes' or 'No', followed by additional reasoning or evidence to support your explanation.} \\
            }
        }
    }
     \\ \\

    \theutterance \stepcounter{utterance}  
    & \multicolumn{4}{p{0.6\linewidth}}{
        \cellcolor[rgb]{0.561,0.737,0.561}{
            \makecell[{{p{\linewidth}}}]{
                \texttt{\tiny{[P1$\rangle$GM]}}
                \texttt{No} \\
\\ 
\texttt{**Reasoning:**} \\
\\ 
\texttt{To answer this question, we must compare the probability of observing \$xevu\$ under two specific counterfactual conditions for an individual who is actually \$rixq\$:} \\
\texttt{1.  **Actual State:** The individual is \$rixq\$.} \\
\texttt{2.  **Counterfactual State:** The individual is *not* \$rixq\$ (while keeping other relevant factors constant, specifically the value of \$zuph\$).} \\
\\ 
\texttt{The question asks: "For those who are \$rixq\$, would it be more likely to see \$xevu\$ if the individual was not \$rixq\$?"} \\
\\ 
\texttt{This requires comparing \$P(xevu \textbackslash{}mid do(rixq=1))\$ with \$P(xevu \textbackslash{}mid do(rixq=0))\$. However, the provided probabilities are conditional on the *observed* values of \$rixq\$ and \$zuph\$, not necessarily interventional probabilities. We must first determine if the causal structure implies confounding or if the observed conditionals suffice to answer the specific counterfactual query based on the "self{-}contained" constraint.} \\
\\ 
\texttt{**Step 1: Analyze the Causal Structure**} \\
\texttt{*   \$Swoy \textbackslash{}rightarrow rixq\$} \\
\texttt{*   \$Swoy \textbackslash{}rightarrow xevu\$ (Direct effect)} \\
\texttt{*   \$rixq \textbackslash{}rightarrow zuph\$} \\
\texttt{*   \$zuph \textbackslash{}rightarrow xevu\$} \\
            }
        }
    }
    & & \\ \\

    \theutterance \stepcounter{utterance}  
    & & & \multicolumn{2}{p{0.3\linewidth}}{
        \cellcolor[rgb]{0.9,0.9,0.9}{
            \makecell[{{p{\linewidth}}}]{
                \texttt{\tiny{[GM$|$GM]}}
                \texttt{No} \\
            }
        }
    }
    & & \\ \\

    \theutterance \stepcounter{utterance}  
    & & & \multicolumn{2}{p{0.3\linewidth}}{
        \cellcolor[rgb]{0.9,0.9,0.9}{
            \makecell[{{p{\linewidth}}}]{
                \texttt{\tiny{[GM$|$GM]}}
                \texttt{no} \\
            }
        }
    }
    & & \\ \\

    \theutterance \stepcounter{utterance}  
    & & & \multicolumn{2}{p{0.3\linewidth}}{
        \cellcolor[rgb]{0.9,0.9,0.9}{
            \makecell[{{p{\linewidth}}}]{
                \texttt{\tiny{[GM$|$GM]}}
                \texttt{game\_result = WIN} \\
            }
        }
    }
    & & \\ \\

\end{supertabular}
}

\end{document}
